\documentclass[11pt,a4paper]{article}

\usepackage[utf8]{inputenc}
\usepackage{amsmath,amssymb}
\usepackage{booktabs}
\usepackage{array}
\usepackage{geometry}
\usepackage{hyperref}
\usepackage{graphicx}

\geometry{margin=2.5cm}

\title{Partial Association Explorer: A Shiny Application for Exploring Conditional Associations in Multivariate Data}

\author{
Thaddée D'haegeleer\textsuperscript{1,*},
Cédric Heuchenne\textsuperscript{1,2},
Antoine Soetewey\textsuperscript{1,2}\\
\textsuperscript{1}TODO: affiliation\\
\textsuperscript{2}TODO: affiliation\\
\textsuperscript{*}Corresponding author: TODO
}

\date{}

\begin{document}

\maketitle

\begin{abstract}
Partial Association Explorer is an open-source R Shiny application for interactive exploration of marginal and conditional associations in multivariate datasets. The software extends the previously published AssociationExplorer application by allowing users to select control variables and assess how pairwise relationships between variables change after adjustment. For numeric--numeric pairs, the app computes partial correlations; for numeric--categorical pairs, it uses an ANCOVA-based partial eta-squared measure; and for categorical--categorical pairs, it introduces a symmetric likelihood-based association measure, denoted $V_L$ in the marginal case and $V_L \mid Z$ in the conditional case. The application also computes statistical significance for global associations and lets users combine effect-size thresholds with $p$-value thresholds when constructing the correlation network. This design helps users distinguish strong and statistically supported associations from patterns that may be induced by confounding variables. The interface is intended for both non-expert users, such as journalists and data communicators, and researchers who need a rapid visual tool for exploratory data analysis. In addition to the network representation, the app provides pairwise diagnostic plots adapted to the variable types, including residual scatter plots, residualized group means, likelihood-based categorical residual tables, Pearson residuals, and optimized submatrix displays for large contingency tables. By integrating conditional association measures into a user-friendly workflow, Partial Association Explorer supports transparent, accessible, and statistically grounded data exploration.
\end{abstract}

\noindent\textbf{Keywords:} R Shiny; exploratory data analysis; partial association; conditional association; correlation network; categorical data

\section*{Metadata}

The metadata associated with the current version of the software is summarized in Table~\ref{tab:metadata}. The application is implemented in R and Shiny and builds on the codebase and workflow of AssociationExplorer \cite{soetewey2026associationexplorer}.

\begin{table}[h]
\centering
\caption{Code metadata.}
\label{tab:metadata}
\begin{tabular}{p{0.08\textwidth}p{0.42\textwidth}p{0.40\textwidth}}
\toprule
Nr. & Code metadata description & Metadata \\
\midrule
C1 & Current code version & TODO: version number \\
C2 & Permanent link to code/repository used for this code version & \url{https://github.com/Thadhaeg/Partial-association-explorer} \\
C3 & Permanent link to reproducible capsule & TODO: Code Ocean, Zenodo, or equivalent archive \\
C4 & Legal code license & TODO: license, e.g. MIT License \\
C5 & Code versioning system used & Git/GitHub \\
C6 & Software code languages, tools, and services used & R, R Shiny \\
C7 & Compilation requirements, operating environments, and dependencies & TODO: R version and package list; current code uses shiny, bslib, dplyr, ggplot2, scales, reactable, tidygraph, visNetwork, readxl, grid, gridExtra, janitor, shinyjs, tibble, shinycssloaders, nnet, and VGAM \\
C8 & Link to developer documentation/manual/screencast & TODO \\
C9 & Support for questions or issues & \url{https://github.com/Thadhaeg/Partial-association-explorer/issues} \\
\bottomrule
\end{tabular}
\end{table}

The remainder of this paper is organized as follows. Section~\ref{sec:motivation} describes the motivation and significance of the software. Section~\ref{sec:software} presents the software architecture and main functionalities. Section~\ref{sec:example} outlines an illustrative workflow. Section~\ref{sec:impact} discusses the expected impact of the application for data exploration, research, and public communication. Section~\ref{sec:conclusion} concludes with limitations and future directions.

\section{Motivation and significance}
\label{sec:motivation}

Exploratory data analysis often begins with a simple question: which variables appear to be related? In datasets that combine numerical, ordinal, and nominal variables, answering this question is not trivial. Standard correlation tools are typically designed for numerical variables, while categorical variables require different association measures and visual summaries. The original AssociationExplorer application addressed this problem by providing an accessible R Shiny interface that automatically selected association measures according to variable type and displayed the results as an interactive correlation network \cite{soetewey2026associationexplorer}.

However, many associations observed in real datasets are not direct relationships between two variables. They may arise because both variables are related to one or more additional variables. This issue is especially important in social surveys, public policy datasets, health datasets, and administrative data, where age, education, income, geography, time, or institutional context can strongly influence many variables simultaneously. A purely marginal analysis can therefore highlight associations that are useful for initial screening but incomplete for interpretation.

Partial Association Explorer was developed to extend the AssociationExplorer workflow from marginal association analysis to conditional association analysis. Users can select one or more control variables and ask whether the relationship between two variables remains after adjustment. This functionality makes it possible to move from a network of raw pairwise associations to a network of residual or adjusted associations. The resulting visualization can help users identify relationships that persist after accounting for plausible confounders, as well as relationships that disappear once those variables are controlled for.

The software is designed for two overlapping audiences. First, it supports non-expert users, including journalists, educators, students, and public data communicators, who need to explore structured datasets without writing code. For this audience, the application provides a guided interface, interpretable plots, and automatic statistical choices. Second, it supports researchers who need a fast exploratory layer before fitting more specialized models. For this audience, the application provides formal test statistics, $p$-values, and model-based adjusted association measures that can guide hypothesis generation, variable selection, and diagnostic work.

The main methodological contribution is the integration of conditional association measures for all pair types into a single interactive workflow. Numeric--numeric relationships are handled through partial correlations. Numeric--categorical relationships are handled through partial eta-squared and extra-sum-of-squares tests. Categorical--categorical relationships are handled through a likelihood-based measure, $V_L$, derived from likelihood-ratio statistics. This measure is symmetric, is naturally connected to likelihood theory, and extends coherently to the conditional case, where the app reports $V_L \mid Z$.

This last point is important because categorical--categorical conditional association is often difficult to present to non-specialist users. Stratified contingency tables can become sparse or impossible to interpret when control variables are numerous or continuous. The application instead uses a regression-based likelihood framework to model the conditional distribution of the two categorical variables given the controls. It then compares a conditional independence model with a model allowing residual association. This gives a global adjusted association measure and local cell diagnostics that can be displayed in pair plots.

\section{Software description}
\label{sec:software}

\subsection{Software architecture}

Partial Association Explorer is implemented as a standalone R Shiny application. The interface is organized into a sequence of tabs corresponding to the main workflow: data upload, variable selection, control-variable selection, correlation network visualization, pairwise plots, and help. The app accepts CSV and Excel files and can optionally read a variable-description file containing human-readable labels for the variables. Variables with no variation are removed during preprocessing because they cannot contribute to association measures or visual summaries.

The application stores the full dataset after preprocessing and lets users select two distinct sets of variables. The first set contains the variables to be visualized in the association network. The second optional set contains control variables used for adjustment. Control variables are not displayed as nodes in the network; instead, they are used when computing conditional associations between the displayed variables. This separation keeps the visual output focused on the relationships of substantive interest while still allowing the analysis to account for background variables.

For every pair of selected variables, the backend determines the variable types and computes an appropriate association measure, a test statistic, and a $p$-value. The app returns three aligned matrices: an association matrix, a matrix describing the type of association measure used, and a $p$-value matrix. These matrices are then filtered according to user-defined thresholds. The correlation network is built from the retained pairs: variables are represented as nodes and associations as edges. Pairwise diagnostic plots are generated for the same retained edges.

\subsection{Association measures and inference}

The application distinguishes between marginal associations, computed without controls, and conditional associations, computed after adjustment for a user-selected set of controls $Z$.

\paragraph{Numeric--numeric pairs.}
For two numeric variables $X$ and $Y$, the marginal association is based on Pearson's correlation coefficient $r_{XY}$, with $R^2_{XY}=r_{XY}^2$ used as the thresholded association strength. When controls are selected, the app computes residuals from regressions of $X$ on $Z$ and $Y$ on $Z$, then computes the correlation between these residuals:
\[
r_{XY\cdot Z} =
\frac{\sum_{i=1}^n \hat{\varepsilon}^X_i \hat{\varepsilon}^Y_i}
{\sqrt{\sum_{i=1}^n(\hat{\varepsilon}^X_i)^2}\sqrt{\sum_{i=1}^n(\hat{\varepsilon}^Y_i)^2}}.
\]
The conditional association strength is $R^2_{XY\cdot Z}=r_{XY\cdot Z}^2$. The corresponding significance test uses the standard partial-correlation statistic with residual degrees of freedom determined by the number of complete observations and controls.

\paragraph{Numeric--categorical pairs.}
For a numeric variable $Y$ and a categorical variable $X$, the marginal association is based on the proportion of variance explained by group membership:
\[
\eta^2 = \frac{\mathrm{SSB}}{\mathrm{SST}},
\]
where $\mathrm{SSB}$ is the between-group sum of squares and $\mathrm{SST}$ is the total sum of squares. In the conditional case, the app fits a full model containing both $X$ and the controls $Z$, and a reduced model containing only $Z$. The partial effect of $X$ is summarized by
\[
\eta^2_{X\mid Z} =
\frac{\mathrm{RSS}_{\mathrm{red}}-\mathrm{RSS}_{\mathrm{full}}}
{\mathrm{RSS}_{\mathrm{red}}-\mathrm{RSS}_{\mathrm{full}}+\mathrm{RSS}_{\mathrm{full}}}.
\]
The associated $p$-value is obtained from the extra-sum-of-squares $F$-test comparing the full and reduced models.

\paragraph{Categorical--categorical pairs.}
For two categorical variables $X$ and $Y$, the app uses a likelihood-based association measure. In the marginal case, it compares the independence model $H_0:X\perp Y$ with the saturated model. Let $\ell_0$ and $\ell_1$ be the maximized log-likelihoods under the null and alternative models, respectively. The likelihood-ratio statistic is
\[
G^2 = 2(\ell_1-\ell_0),
\]
and the association measure is
\[
V_L = \sqrt{1-\exp(-G^2/n)}.
\]
This quantity is a monotone transformation of the Cox--Snell pseudo-$R^2$ and can be interpreted as a likelihood-based analogue of a correlation coefficient for categorical variables. Unlike Cramer's $V$, it is derived directly from the likelihood-ratio statistic used for inference, which makes it convenient for extension to conditional models.

In the conditional case, the app models the joint conditional distribution of $(X,Y)$ given $Z$ using a structured multinomial framework. The null model imposes conditional independence,
\[
\Pr(X=i,Y=j\mid Z=z)=\Pr(X=i\mid Z=z)\Pr(Y=j\mid Z=z),
\]
while the alternative model allows residual interaction terms between categories of $X$ and $Y$. The likelihood-ratio statistic compares these nested models, and the app reports the adjusted measure $V_L\mid Z$:
\[
V_L\mid Z = \sqrt{1-\exp(-G^2_Z/n)}.
\]
The corresponding $p$-value is computed from the asymptotic chi-square distribution with degrees of freedom equal to the number of residual interaction parameters.

\subsection{Software functionalities}

The major functionalities of Partial Association Explorer include:

\begin{itemize}
\item \textbf{Data upload and preprocessing.} The app supports CSV and Excel datasets. It removes variables with only one unique value and can integrate an optional variable-description file to improve interpretability.
\item \textbf{Variable and control selection.} Users select variables to display in the network and, separately, optional control variables used to compute adjusted associations.
\item \textbf{Dynamic association filtering.} Users set effect-size thresholds for numeric/numeric and numeric/categorical associations, a separate threshold for categorical/categorical associations, and a significance threshold based on $p$-values. The network and pair plots update reactively.
\item \textbf{Interactive correlation network.} The app displays retained associations as edges between variable nodes. Edge properties encode association strength and type, helping users identify relationships that remain after adjustment.
\item \textbf{Pairwise diagnostic plots.} Numeric--numeric pairs are shown with scatter plots or residual scatter plots. Numeric--categorical pairs are shown with group means or residualized means. Categorical--categorical pairs are shown using observed-minus-expected deviations and Pearson residuals, with local diagnostics derived from the null model.
\item \textbf{Large contingency-table handling.} For categorical--categorical pair plots involving many categories, the app identifies an informative submatrix so that the displayed table remains readable and focused on the strongest local departures.
\item \textbf{Axis reversal for scatter plots.} Users can reverse scatter plot axes, which improves interpretability when the relationship is easier to read in the opposite orientation.
\end{itemize}

\subsection{Sample code snippet}

The following representative excerpt illustrates the core logic used by the application: detect the pair type, compute a marginal or conditional association depending on whether controls were selected, and store both effect sizes and $p$-values for filtering.

\begin{verbatim}
if (is_num1 && is_num2) {
  if (has_controls) {
    resid_x <- partial_residuals(x, control_data)
    resid_y <- partial_residuals(y, control_data)
    r <- cor(resid_x, resid_y, use = "complete.obs")
    p_value <- p_value_partial_cor(r, n_eff, k_controls)
    association <- r
    association_type <- "Partial r"
  } else {
    r <- cor(x, y, use = "complete.obs")
    association <- r
    association_type <- "Pearson's r"
  }
} else if (is_num1 != is_num2) {
  result <- calculate_partial_eta_squared_with_F(
    num_var, cat_var, control_data
  )
  association <- result$eta
  p_value <- result$p_value
  association_type <- result$type
} else {
  result <- if (has_controls) {
    compute_conditional(x, y, control_data)
  } else {
    compute_unconditional(x, y)
  }
  association <- result$VL
  p_value <- result$p_value
  association_type <- if (has_controls) "VL|Z" else "VL"
}
\end{verbatim}

This structure preserves the accessible workflow of AssociationExplorer while adding conditional statistical modelling behind the interface.

\section{Illustrative example}
\label{sec:example}

TODO: Insert the final dataset and figures used for the submission. A natural illustrative example would reuse the European Social Survey workflow from the first AssociationExplorer paper and add a control-variable scenario. For example, users could explore associations between political attitudes, media consumption, trust, health, and demographic variables among Belgian respondents, first without controls and then after controlling for age, gender, education, income, or country/region where appropriate.

The example workflow would proceed as follows. First, the user uploads the dataset and, optionally, a description file containing variable labels from the survey codebook. Second, the user selects the substantive variables to display in the association network. Third, the user selects one or more control variables. Fourth, the app computes all pairwise adjusted associations among the displayed variables and filters them by both effect-size and significance thresholds. Finally, the user opens the pair-plot tab to inspect the retained associations locally.

Such an example would show the added value of the new software compared with the original application. A marginal network may reveal a dense set of associations, some of which reflect common dependence on demographic or contextual variables. After selecting controls, the adjusted network may become sparser and highlight relationships that persist after accounting for those variables. The pair plots then explain the retained relationships: residual scatter plots for numeric pairs, residualized group means for numeric--categorical pairs, and cell-wise residual tables for categorical--categorical pairs.

Figures to include:

\begin{itemize}
\item Fig.~1: Data upload tab.
\item Fig.~2: Variable and control-variable selection tab.
\item Fig.~3: Marginal association network.
\item Fig.~4: Conditional association network after selecting controls.
\item Fig.~5: Residual scatter plot for a numeric--numeric pair.
\item Fig.~6: Residualized means plot for a numeric--categorical pair.
\item Fig.~7: Pearson residual/deviation table for a categorical--categorical pair.
\end{itemize}

\section{Impact}
\label{sec:impact}

Partial Association Explorer expands the scope of user-friendly association exploration from marginal relationships to adjusted relationships. This is a meaningful practical step because many users can identify correlations visually, but fewer have access to tools that make conditional association analysis accessible without requiring code. By integrating control-variable selection directly into the workflow, the app helps users ask more careful exploratory questions: does this association remain after accounting for background variables? Which relationships are robust to adjustment? Which apparent links are likely explained by a shared dependence on other variables?

For journalists and public communicators, the application can support responsible data storytelling. Open datasets often contain many variables and many possible pairwise relationships. Without a guided interface, it is easy either to miss potentially important patterns or to overinterpret unadjusted associations. The combination of effect-size thresholds, $p$-value filtering, and interpretable pair plots provides a transparent way to screen data and identify relationships that deserve further investigation. The app does not replace domain expertise or causal inference, but it can make the first exploratory stage more systematic and more accountable.

For researchers, the software can serve as a rapid diagnostic tool before formal modelling. The conditional network may suggest candidate relationships, reveal redundancy among variables, highlight confounding patterns, or indicate where categorical interactions deserve further analysis. The likelihood-based $V_L$ and $V_L\mid Z$ measures are especially useful for datasets containing nominal variables, where many standard correlation tools are inadequate or require transformations that obscure interpretation.

The app also has pedagogical value. It allows students to compare marginal and conditional association in the same dataset and to see how adjustment changes both the global network and the local plots. This can help teach central statistical ideas such as confounding, residualization, partial correlation, analysis of covariance, likelihood-ratio tests, and conditional independence using interactive examples rather than isolated formulas.

Compared with the original AssociationExplorer application, the updated version implements several planned extensions: partial correlations, axis reversal for scatter plots, significance-based filtering, and richer categorical diagnostics. It therefore represents not only a technical update but a conceptual extension from descriptive bivariate exploration toward model-assisted exploratory analysis.

\section{Conclusions}
\label{sec:conclusion}

Partial Association Explorer is an open-source R Shiny application for exploring marginal and conditional associations in multivariate datasets. It extends the AssociationExplorer workflow by adding control-variable selection, partial association measures, likelihood-based categorical association measures, significance filtering, and pairwise plots designed for adjusted interpretation. The app remains accessible to non-technical users while offering researchers a useful exploratory tool for mixed-type datasets.

The current version should be interpreted as an exploratory tool rather than a causal inference platform. Conditional associations depend on the selected controls and on the adequacy of the underlying regression models. In particular, the categorical--categorical conditional association measure relies on model-based expected counts under conditional independence. Small samples, sparse category combinations, highly imbalanced tables, and model misspecification may affect the reliability of both the global statistics and local residual diagnostics.

Future development will focus on documenting these assumptions clearly, improving computational efficiency for high-dimensional datasets, adding further export and reporting options, and validating the workflow through applied case studies. Additional extensions may include survey weights, more flexible handling of nonlinear control effects, user guidance for interpreting adjusted associations, and richer comparison tools between marginal and conditional networks.

\section*{CRediT authorship contribution statement}

Thaddée D'haegeleer: TODO. Cédric Heuchenne: TODO. Antoine Soetewey: TODO.

\section*{Declaration of competing interest}

The authors declare that they have no known competing financial interests or personal relationships that could have appeared to influence the work reported in this paper.

\section*{Acknowledgements}

TODO: Add funding, project acknowledgements, institutional acknowledgements, and thanks to contributors or testers.

\begin{thebibliography}{99}

\bibitem{soetewey2026associationexplorer}
A. Soetewey, C. Heuchenne, A. Claes, A. Descampe,
AssociationExplorer: A user-friendly Shiny application for exploring associations and visual patterns,
SoftwareX 33 (2026) 102483.
\url{https://doi.org/10.1016/j.softx.2025.102483}

\bibitem{r}
R Core Team,
R: A language and environment for statistical computing,
R Foundation for Statistical Computing, Vienna, Austria, 2024.
\url{https://www.R-project.org/}

\bibitem{shiny}
W. Chang, J. Cheng, J.J. Allaire, C. Sievert, B. Schloerke, Y. Xie, J. Allen, J. McPherson, A. Dipert, B. Borges,
shiny: Web application framework for R,
R package, 2024.
\url{https://CRAN.R-project.org/package=shiny}

\bibitem{visnetwork}
Almende B.V. and contributors, B. Thieurmel,
visNetwork: Network visualization using vis.js library,
R package, 2022.
\url{https://CRAN.R-project.org/package=visNetwork}

\bibitem{vglm}
T.W. Yee,
The VGAM package for categorical data analysis,
Journal of Statistical Software 32 (10) (2010) 1--34.
\url{https://doi.org/10.18637/jss.v032.i10}

\end{thebibliography}

\end{document}
